% Section 1 -- introduction.  v3, 13 Oct 2026 (FINAL_SUBMISSION_REVIEW_ASTRA: compressed; three contributions; residual-polynomial
% wording; no baseline-superiority claim).  Numbers: manuscript/tables/numbers.txt.
\section{Introduction}\label{sec:intro}
Many tasks in computational mechanics solve the same elliptic problem for many coefficient fields: uncertainty
quantification for random media, inverse problems and design loops over heterogeneous materials. Neural operators
\cite{kovachki2023,li2021fno,tripura2023wno} promise to amortise such solves by learning the map from coefficient to
solution. At high contrast this is difficult. When the coefficient varies by four orders of magnitude or more, a small nodal
error can hide a large error in the energy norm, the norm in which the discrete problem is posed and in which multigrid and
coarse-space theory measure accuracy. Learned surrogates are moreover often compared with classical methods that receive
less information or are not run to a matched accuracy, which overstates their benefit \cite{mcgreivy2024,chatain2025}.

Classical solvers treat high contrast with operator-dependent multilevel structure. Algebraic multigrid builds its
interpolation from the matrix \cite{vanek1996,xu2017acta}; operator-adapted wavelets
\cite{amaratunga2006,sudarshan2006,sudarshan2005thesis,dheedene2005thesis} and gamblets \cite{owhadi2017} decompose the
solution space into levels that are energy-orthogonal for the given operator; localized orthogonal decomposition and
spectral coarse spaces make the coarse level robust to the contrast
\cite{malqvist2014,peterseim2016,hellman2017,galvis2010,spillane2014geneo}. Learned multigrid methods keep this structure
and learn some of its components, but they iterate the resulting cycle or use it as a preconditioner
\cite{greenfeld2019,luz2020,huang2023,wangluo2026,zhanglisaad2026}. We ask a narrower, architectural question. Given one
pass through a fixed multilevel architecture on a hierarchy adapted to the operator of each sample, can learned smoothers
make that pass measurably more accurate than the same pass with classical smoothers, when both use the same hierarchy,
coarse solution and correction ingredients? How does that answer depend on the classical smoothing budget, and what does the
learned accuracy cost compared with classical solvers run to the same accuracy?

To answer it we construct OA-VCB, a learned V-cycle on an unstructured finite element hierarchy
(Section~\ref{sec:method}). Its transfers are computed from the operator of each sample by a few Jacobi steps towards the
ideal, Schur-complement prolongation of operator-adapted wavelets; its coarse operators are Galerkin products; its coarsest
level is solved exactly. The only learned parts are the smoothers, small nonlinear networks that see the residual and the
local operator weights. Floating conductive regions, whose indicator functions are near-kernel vectors of the operator, are
detected from the matrix and deflated, and the output is the Galerkin projection of the cycle's directions. Every learned
result is paired, sample by sample, with a \emph{classical reference}: the same cycle with degree-four Chebyshev smoothers
and the same coarse solution, deflation and output projection.

Our contributions are the following.
\begin{itemize}
\item A learned V-cycle for unstructured finite element discretisations of high-contrast elliptic problems whose transfers,
coarse operators and deflation vectors are computed from each operator. With $\nu$ update stages per smoother it represents
every classical cycle on the same hierarchy whose smoothers have a fixed, level-independent residual polynomial of degree at
most $\nu$, including weighted Jacobi and fixed-interval Chebyshev smoothing; it is positively homogeneous in the load; and
its output is never worse in the energy norm than the coarse approximation it starts from (Section~\ref{sec:method}).
\item A controlled paired evaluation on three held-out data seeds. With four update stages per pre- and post-smoother, one
learned pass is about ten times more accurate in the energy norm than the degree-four Chebyshev reference on conductive
channels ($G = 9.7$--$10.2$, more accurate on every test sample; $9.7$--$10.7$ over three training initialisations) and about
twice as accurate on insulating barriers (Section~\ref{sec:accuracy}). The margin depends strongly on the classical smoothing
budget: sixteen Chebyshev steps narrow it to $1.8$ on channels and reverse the paired advantage on barriers, and three
classical passes are more accurate than the learned pass on both families at about a third of its time
(Section~\ref{sec:avt}). Ablations show that the
operator-adapted transfers, the near-kernel deflation, the Galerkin output and the nonlinearity of the smoothers each
contribute under this training protocol (Section~\ref{sec:design}).
\item A measured account of the limits (Sections~\ref{sec:cost} and~\ref{sec:failure}): a fixed-load network loses its
advantage on other loads, while distributional training restores it on that distribution; substantial error remains on
floating islands, which a harmonic enrichment with a prespecified written protocol reduces only modestly; and on one CPU the
implementation is one to two orders of magnitude more expensive than converged-to-target classical solvers at matched
accuracy, with no break-even on the tested configurations.
\end{itemize}
Our objective is to isolate the contribution of learned smoothing within an operator-adapted multilevel construction. The
matched classical reference is a control for this architectural question, not a claim of equal work; practical solver
comparisons are made separately, at matched accuracy, and give no computational reason to replace the tested solvers with
OA-VCB. The structural properties are proved; the accuracy results are empirical, and we give no convergence or
approximation theory for the learned smoothers. Learned operator baselines of similar size, adapted to this setting, are
reported for context in \appref{app:baselines}; their rows assess our adaptations, not the published architectures on their
own benchmarks. Section~\ref{sec:related} reviews related work and Section~\ref{sec:discussion} draws the implications of the results
for the design and evaluation of learned multilevel components.
